%==============================================================
\chapter{Molecule Data}
\label{app:molecule_data}
%==============================================================
% ------------------------------- OCS -------------------------------
\begin{table}[htbp]
  \centering
  \caption{Rotational parameters of OCS in the gas phase and inside
           \textsuperscript{4}He nanodroplets.}
  \label{tab:ocs-params}
  \begin{tabular}{|*{6}{c|}}
    \hline
    % --- 2 cells ---
    \multicolumn{3}{|c|}{\textbf{Gas phase}} &
    \multicolumn{3}{c|}{\textbf{Helium environment}} \\
    \hline
    % --- 3 cells ---
    \multicolumn{2}{|P{\spanII}|}{%
      \pcell{$B_\mathrm{gas}$~\cite{lovas1974}}{$6.0815\,$GHz}} &
    \multicolumn{2}{P{\spanII}|}{%
      \pcell{$B_\mathrm{gas}/B_\mathrm{He}$}{$2.78$}} &
    \multicolumn{2}{P{\spanII}|}{%
      \pcell{$B_\mathrm{He}$~\cite{chatterleyRotationalCoherenceSpectroscopy2020}}{$2.19\,$GHz}} \\
    \hline
    % --- 2 cells ---
    \multicolumn{3}{|P{\spanIII}|}{%
      \pcell{$D_\mathrm{gas}$~\cite{lovas1974}}{$1.30\,$kHz}} &
    \multicolumn{3}{P{\spanIII}|}{%
      \pcell{$D_\mathrm{He}$~\cite{chatterleyRotationalCoherenceSpectroscopy2020}}{$9.5\,$MHz}} \\
    \hline
    % --- 1 cell ---
    \multicolumn{6}{|P{\spanVI}|}{%
      \pcell{$\Delta\alpha$~\cite{bridge1966}}{%
        $4.04\,\text{\AA}^3$}} \\
    \hline
    % --- 2 cells ---
    \multicolumn{3}{|P{\spanIII}|}{%
      \pcell{$\nu_\mathrm{gas}(J\!=\!0\!\to\!2)$}{$36.5\,$GHz}} &
    \multicolumn{3}{P{\spanIII}|}{%
      \pcell{$\nu_\mathrm{He}(L\!=\!0\!\to\!2)$}{$13.1\,$GHz}} \\
    \hline
  \end{tabular}
\end{table}

% ------------------------------- CS2 -------------------------------
\begin{table}[htbp]
  \centering
  \caption{Rotational parameters of CS\textsubscript{2} in the gas phase and inside
           \textsuperscript{4}He nanodroplets. Only even $J$ ($L$) are populated.}
  \label{tab:cs2-params}
  \begin{tabular}{|*{6}{c|}}
    \hline
    \multicolumn{3}{|c|}{\textbf{Gas phase}} &
    \multicolumn{3}{c|}{\textbf{Helium environment}} \\
    \hline
    \multicolumn{2}{|P{\spanII}|}{%
      \pcell{$B_\mathrm{gas}$~\cite{herzberg1966}}{$3.2715\,$GHz}} &
    \multicolumn{2}{P{\spanII}|}{%
      \pcell{$B_\mathrm{gas}/B_\mathrm{He}$}{$4.5$}} &
    \multicolumn{2}{P{\spanII}|}{%
      \pcell{$B_\mathrm{He}$~\cite{chatterleyRotationalCoherenceSpectroscopy2020}}{$0.72\,$GHz}} \\
    \hline
    \multicolumn{3}{|P{\spanIII}|}{%
      \pcell{$D_\mathrm{gas}$~\cite{herzberg1966}}{$\approx 0.4\,$kHz}} &
    \multicolumn{3}{P{\spanIII}|}{%
      \pcell{$D_\mathrm{He}$~\cite{chatterleyRotationalCoherenceSpectroscopy2020}}{$0.9\,$MHz}} \\
    \hline
    \multicolumn{6}{|P{\spanVI}|}{%
      \pcell{$\Delta\alpha$~\cite{bridge1966}}{%
        $\approx 9.6\,\text{\AA}^3$}} \\
    \hline
    \multicolumn{3}{|P{\spanIII}|}{%
      \pcell{$\nu_\mathrm{gas}(J\!=\!0\!\to\!2)$}{$19.6\,$GHz}} &
    \multicolumn{3}{P{\spanIII}|}{%
      \pcell{$\nu_\mathrm{He}(L\!=\!0\!\to\!2)$}{$4.3\,$GHz}} \\
    \hline
  \end{tabular}
\end{table}
% Supplementary derivations, extended data tables, additional figures.
\todo{Find citations and add ratio of D}

\begin{table}[htbp]
  \centering
\caption{Rotational parameters of the nitric oxide dimer, \textit{cis}-ONNO,
in the gas phase and inside \textsuperscript{4}He nanodroplets.}
\label{tab:no2-params}
\begin{tabular}{|*{6}{c|}}
\hline
% --- 2 cells ---
\multicolumn{3}{|c|}{\textbf{Gas phase}} &
\multicolumn{3}{c|}{\textbf{Helium environment}} \\
\hline
% --- 3 cells ---
\multicolumn{2}{|P{\spanII}|}{%
\pcell{$B_{yz}^\mathrm{gas}$~\cite{howard1993}}{$5.10\,$GHz}} &
\multicolumn{2}{P{\spanII}|}{%
\pcell{$B_{yz}^\mathrm{gas}/B_{yz}^\mathrm{He}$}{$1.9$}} &
\multicolumn{2}{P{\spanII}|}{%
\pcell{$B_{yz}^\mathrm{He}$~\cite{macphail-bartleyControlMolecularRotation2026a}}{$2.76(6)\,$GHz}} \\
\hline
% --- 2 cells ---
\multicolumn{3}{|P{\spanIII}|}{%
\pcell{$D_\mathrm{gas}$~\cite{howard1993}}{%
$\approx 1\times10^{-6}\,$cm\textsuperscript{-1}}} &
\multicolumn{3}{P{\spanIII}|}{%
\pcell{$D_\mathrm{He}$}{---}} \\
\hline
% --- 1 cell ---
\multicolumn{6}{|P{\spanVI}|}{%
\pcell{$\Delta\alpha$~\cite{cccbdb} (calculated)}{%
$3.4\,\text{\AA}^3$}} \\
\hline
% --- 2 cells ---
\multicolumn{3}{|P{\spanIII}|}{%
\pcell{$\nu_\mathrm{gas}(J\!=\!0\!\to\!2)$}{$31\,$GHz}} &
\multicolumn{3}{P{\spanIII}|}{%
\pcell{$\nu_\mathrm{He}(J\!=\!0\!\to\!2)$}{$16.8(4)\,$GHz}} \\
\hline
\end{tabular}
\end{table}

\chapter{Different molecules/ experiments}


\chapter{Useful information when conducting experiments}
\section{Chamber Pressures}
\subsection{Vacuum setup}
\section{Optical alignment}
\section{Switching between Jet and Droplets}
\section{Finding the correct TOF window}


\chapter{AI-generated content}
\label{sec:ai}
Content wise, artificial intelligence was only used in the theoretical background section \ref{sec:theory}, with the purpose of accelerating and enhancing the explanation of already known science and math. Every result, if not marked different, was checked and verified by me. Places where AI altered results are used, have footnote markings. 

\section{Numerical illustration of rotational ladder climbing}
\label{app:centrifuge-tdse}
For the purpose of displaying the rotational ladder climbing mechanism, I used Claude to write a script that simulates the dynamics of the rotating frame Hamiltonian \eqnref{eq:rotating_frame_hamiltonian}. This code is \textbf{not} proofread by me, but sanity checked through the outcomes of the simulation.  It implements a direct
numerical solution of the time-dependent Schr\"odinger equation (TDSE) for a rigid rotor driven by an optical centrifuge. The intention of this calculation is of pedagogical nature: it is used to generate the figures that show how population is transferred stepwise up the rotational ladder, and it is not presented as a quantitative model of the experiment reported in this work.

The following explanation of the code is generated by AI.

\subsection{Physical model}
\label{app:centrifuge-model}

The optical centrifuge is formed from two counter-rotating, counter-chirped circularly polarised beams, whose superposition is a linearly polarised field
whose polarisation direction $\phi_d(t)$ rotates with increasing angular
velocity. After averaging over the optical carrier, a linear molecule of
polarisability components $\alpha_\parallel,\alpha_\perp$ experiences the
interaction
%
\begin{equation}
  V(t) = -\frac{\varepsilon_0^2(t)}{4}
         \left[\alpha_\perp
               + \Delta\alpha\,\sin^2\theta\,\cos^2\!\bigl(\varphi-\phi_d(t)\bigr)\right],
  \qquad \Delta\alpha = \alpha_\parallel-\alpha_\perp ,
  \label{eq:app-lab-interaction}
\end{equation}
%
where $\theta,\varphi$ are the polar and azimuthal angles of the molecular
axis and $\varepsilon_0(t)$ is the field envelope. Transforming to the frame
co-rotating with the polarisation by means of the unitary operator
$\hat U(t)=\exp[-\mathrm{i}\phi_d(t)\hat J_z]$ removes the explicit
$\phi_d$-dependence at the cost of an inertial term, giving the working
Hamiltonian
%
\begin{equation}
  \hat{\tilde H}(t)
  = B\hat J^{2}
    - \hbar\Omega(t)\hat J_{z}
    - \frac{\varepsilon_0^2(t)}{4}
      \bigl[\alpha_\perp + \Delta\alpha\,\sin^2\theta\cos^2\varphi\bigr],
  \qquad \Omega=\dot\Phi .
  \label{eq:app-rotating-frame}
\end{equation}
%
This change of frame is exact; no approximation is involved, and no
restriction is placed on $\dot\Omega$. Two further points about
Eq.~\eqref{eq:app-rotating-frame} as implemented:

\begin{enumerate}
  \item The isotropic term $-\tfrac{1}{4}\varepsilon_0^2\alpha_\perp$ is
        proportional to the identity operator. It therefore contributes only a
        time-dependent global phase and cannot affect any observable. It is
        omitted from the computation.
  \item Writing $U(t)=\Delta\alpha\,\varepsilon_0^2(t)/4$ and
        $\Omega(t)=2\pi f_d(t)$, the propagated Hamiltonian in frequency units
        is
        \begin{equation}
          \tilde H(t)/h = B\,J(J+1) - f_d(t)\,M - \frac{U(t)}{h}\,Q_x,
          \qquad Q_x \equiv \sin^2\theta\cos^2\varphi .
          \label{eq:app-working-hamiltonian}
        \end{equation}
\end{enumerate}

\paragraph{Basis.}
The wavefunction is expanded in field-free rigid-rotor eigenstates
$\ket{J,M}$. Because $Q_x$ is a rank-2 spherical tensor with components
$q=0,\pm2$, it obeys the selection rules $\Delta J=0,\pm2$ and
$\Delta M=0,\pm2$ and therefore conserves the parity of both $J$ and $M$.
Starting from $\ket{0,0}$, the accessible Hilbert space is exactly the set of
states with even $J$ and even $M$; restriction to this sector is a rigorous
symmetry reduction rather than an approximation. The basis is truncated at
$J_{\max}$, giving $81$ states for $J_{\max}=16$.

The interaction matrix is built from the spherical-harmonic expansion
%
\begin{equation}
  Q_x = \frac{1}{3}
        - \sqrt{\frac{4\pi}{45}}\,Y_2^{0}
        + \sqrt{\frac{2\pi}{15}}\left(Y_2^{2}+Y_2^{-2}\right),
  \label{eq:app-qx-expansion}
\end{equation}
%
with matrix elements evaluated in closed form from Wigner $3j$ symbols,
%
\begin{equation}
  \begin{aligned}
  \braket{J'M'|Y_k^q|JM}
  \\= (-1)^{M'}\sqrt{\frac{(2J'+1)(2k+1)(2J+1)}{4\pi}}
    \begin{pmatrix} J' & k & J\\ 0&0&0\end{pmatrix}
    \begin{pmatrix} J' & k & J\\ -M' & q & M\end{pmatrix}.
  \end{aligned}
  \label{eq:app-wigner}
\end{equation}
%
Both the $q=+2$ and $q=-2$ terms are retained. The calculation therefore makes
\emph{no} rotating-wave approximation and imposes \emph{no} constraint
$M=J$; couplings that violate the resonant selection rule $\Delta J=\Delta M$
are present, as are all $\Delta J=0$ (dressing) terms.

\paragraph{Pulse and chirp schedule.}
The laser intensity envelope is Gaussian with full width at half maximum
$\tau_{\mathrm{FWHM}}$, centred at $t=0$; since the polarisability interaction
is linear in intensity,
%
\begin{equation}
  U(t)/h = (U_0/h)\,\exp\!\left[-4\ln 2\,(t/\tau_{\mathrm{FWHM}})^2\right].
  \label{eq:app-envelope}
\end{equation}
%
The rotation frequency is held at zero, ramped linearly across the FWHM
interval, and then held constant:
%
\begin{equation}
  f_d(t) = f_{\mathrm{final}}\times
  \begin{cases}
    0, & t \le -\tau_{\mathrm{FWHM}}/2,\\[2pt]
    \bigl(t+\tau_{\mathrm{FWHM}}/2\bigr)/\tau_{\mathrm{FWHM}},
      & |t| < \tau_{\mathrm{FWHM}}/2,\\[2pt]
    1, & t \ge +\tau_{\mathrm{FWHM}}/2 .
  \end{cases}
  \label{eq:app-chirp}
\end{equation}
%
The field is thus already at half its peak value when the rotation begins, so
the molecule is adiabatically dressed into a pendular state before the
centrifuge starts, and the field is still at half peak when the rotation stops,
so the final state is released by a slow turn-off at constant $\Omega$. Each
value of $f_{\mathrm{final}}$ corresponds to a different constant angular
acceleration.

Propagation begins and ends where $U/U_0=10^{-4}$ rather than at the edges of
the plotted window, so that the initial condition $\ket{0,0}$ is imposed in an
essentially field-free region and the reported final populations are
field-free populations rather than populations of dressed states. Because the
frame transformation is diagonal in $M$ and does not mix $J$, the quantities
$P_J$ and $\langle J\rangle$ are identical in the rotating and laboratory
frames.

\subsection{Numerical method and parameters}
\label{app:centrifuge-numerics}

The TDSE is integrated with a symmetric (Strang) split-operator propagator,
%
\begin{equation}
  \exp\!\left[-\tfrac{\mathrm{i}}{\hbar}\tilde H\,\delta t\right]
  \simeq
  \mathrm{e}^{-\mathrm{i}\tilde H_{\mathrm{diag}}\delta t/2\hbar}\,
  \mathrm{e}^{-\mathrm{i}\tilde H_{\mathrm{int}}\delta t/\hbar}\,
  \mathrm{e}^{-\mathrm{i}\tilde H_{\mathrm{diag}}\delta t/2\hbar}
  + \mathcal{O}(\delta t^{3}),
  \label{eq:app-strang}
\end{equation}
%
where $\tilde H_{\mathrm{diag}}/h = B J(J+1) - f_d(t)M$ is diagonal in
$\ket{J,M}$ and $\tilde H_{\mathrm{int}}/h = -(U(t)/h)\,Q_x$ is diagonal in the
eigenbasis of $Q_x$. Since $Q_x$ is time-independent, it is diagonalised once
and reused. Each factor is separately unitary. All values of
$f_{\mathrm{final}}$ are propagated simultaneously as columns of a single
matrix; this is mathematically identical to independent runs.

Parameters are collected in Table~\ref{tab:app-parameters}.

\begin{table}[htbp]
  \centering
  \caption{Parameters of the ladder-climbing calculation.}
  \label{tab:app-parameters}
  \begin{tabular}{llll}
    \toprule
    Quantity & Symbol & Value & Source / note \\
    \midrule
    Rotational constant & $B$ & \SI{6.081492}{\giga\hertz}
      & gas-phase $B_0$ of $^{16}$O$^{12}$C$^{32}$S \\
    Centrifugal distortion & $D$ & --- & omitted \\
    Polarisability anisotropy & $\Delta\alpha$ & \SI{5.34}{\cubic\angstrom}
      & static value \\
    Peak intensity & $I_0$ & \SI{1e11}{\watt\per\square\centi\meter} & \\
    Peak coupling & $U_0/h$ & \SI{168.9}{\giga\hertz} & derived \\
    Dimensionless coupling & $U_0/B$ & $27.8$ & derived \\
    Intensity FWHM & $\tau_{\mathrm{FWHM}}$ & \SI{320}{\pico\second} & \\
    Final frequency range & $f_{\mathrm{final}}$ & $0$--\SI{120}{\giga\hertz}
      & 200 values \\
    Basis truncation & $J_{\max}$ & $16$ & 81 states \\
    Integration step & $\delta t$ & \SI{0.05}{\pico\second} & \\
    Initial state & --- & $\ket{J=0,M=0}$ & $T=\SI{0}{\kelvin}$ \\
    \bottomrule
  \end{tabular}
\end{table}

\subsection{Verification}
\label{app:centrifuge-verification}

The following checks were performed and are reported here so that the
calculation can be reproduced or refuted.

\begin{enumerate}
  \item \textbf{Operator construction.} The diagonal element
        $\braket{0,0|Q_x|0,0}$ reproduces the analytic spherical average
        $\langle\sin^2\theta\cos^2\varphi\rangle = 1/3$ to machine precision,
        and all eigenvalues of the truncated $Q_x$ matrix lie within the range
        $[0,1]$ of the exact operator.
  \item \textbf{Time-step convergence.} Reducing $\delta t$ by a factor of
        four changes the final $\langle J\rangle$ in the seventh decimal
        place. Note that norm conservation is \emph{not} a useful convergence
        diagnostic here, because the split propagator is unitary by
        construction; the norm is conserved to roundoff regardless of whether
        $\delta t$ is adequate.
  \item \textbf{Basis convergence.} Increasing $J_{\max}$ from 16 to 24 leaves
        the final populations unchanged to six decimal places. The population
        residing in the highest retained $J$ manifold remains below $10^{-12}$
        throughout.
  \item \textbf{Consistency with the analytic ladder picture.} The final
        rotational state agrees with the number of two-photon resonances
        $f_d = B(2J+3)$ traversed during the chirp, and simultaneously with the
        classical trapped-rotor estimate $J\approx f_d/2B$; at these parameters
        the two predictions coincide. Representative values are given in
        Table~\ref{tab:app-results}.
  \item \textbf{Role of non-resonant couplings.} At these parameters
        $99.99\,\%$ of the final population lies in the chain
        $C\equiv J-M=0$, and imposing the rotating-wave approximation by hand
        changes the final populations by $<10^{-3}$. The non-resonant terms are
        thus retained but numerically inactive here; their inclusion matters
        for weaker fields or thermal initial conditions.
\end{enumerate}

\begin{table}[htbp]
  \centering
  \caption{Representative results. $\sigma_J$ is the standard deviation of the
           final rotational distribution.}
  \label{tab:app-results}
  \begin{tabular}{S[table-format=3.0] S[table-format=2.3] S[table-format=1.2] S[table-format=1.3]}
    \toprule
    {$f_{\mathrm{final}}$ (GHz)} & {$\langle J\rangle_{\mathrm{final}}$}
      & {$\sigma_J$} & {$P_{J_{\mathrm{dom}}}$} \\
    \midrule
     60 &  4.002 & 0.07 & 0.999 \\
    100 &  7.991 & 0.14 & 0.995 \\
    120 &  9.986 & 0.18 & 0.993 \\
    \bottomrule
  \end{tabular}
\end{table}

At the parameters of Table~\ref{tab:app-parameters} the transfer is close to
adiabatic: more than $99\,\%$ of the population ends in a single rotational
state with $\sigma_J\lesssim0.2$. In the dimensionless parametrisation of
Armon and Friedland \cite{armonQuantumClassicalDynamics2017}, the calculation sits near
the boundary between quantum ladder climbing and classical autoresonance
($P_1\approx15$, $P_2\approx1.1$ at the largest final frequency), with the
autoresonant capture criterion $P_1P_2>1/2$ satisfied by more than an order of
magnitude. Weaker fields or faster chirps move the system towards the
non-adiabatic regime, in which population is left behind at successive
crossings and the final distribution broadens.

\subsection{Limitations}
\label{app:centrifuge-limitations}

The following are omitted from the model and bound its interpretation.

\begin{itemize}
  \item \textbf{Centrifugal distortion.} The rotor is treated as rigid. The
        correction to the two-photon resonance grows as $J^{3}$ and is
        negligible over the range of $J$ shown here, but the model must not be
        extrapolated to superrotor states without reinstating $D$.
  \item \textbf{Temperature.} The calculation starts from the single pure
        state $\ket{0,0}$. A thermal ensemble populates many initial
        $\ket{J_0,M_0}$, which are captured into the rotating well with
        differing efficiency; the resulting bifurcation between trapped and
        untrapped populations is the dominant source of the broad final
        distributions seen experimentally, and it is absent here by
        construction.
  \item \textbf{Polarisability.} The static $\Delta\alpha$ is used in place of
        the frequency-dependent value at the laser wavelength.
  \item \textbf{Experimental averaging.} No focal-volume or intensity
        averaging is performed; a single peak intensity is assumed.
  \item \textbf{Dissipation.} The propagation is strictly unitary. Collisional
        and environmental decoherence are not included.
\end{itemize}